\documentclass[a4paper, 12pt]{report}

\author{Deniz Güven}

\usepackage[ngerman]{babel}
\usepackage[utf8]{inputenc}
\usepackage[T1]{fontenc}
\usepackage{minted}
\usepackage{tikz}
\usetikzlibrary{arrows, automata}
\usepackage{xcolor}
\usepackage{forest}
%\usepackage[margin=23mm, top=40mm, headheight=24pt, bottom=30mm]{geometry}
\usepackage{graphicx}
\usepackage{fancyhdr}
\usepackage{amsmath, amsfonts, amssymb, amsthm}
\usepackage{mathtools}
\usepackage{stmaryrd}
\usepackage[parfill]{parskip}
\usepackage[colorlinks]{hyperref}
\usepackage[]{cleveref} % to be loaded after hyperref
\crefname{equation}{}{}

\begin{document}

\begin{titlepage}
  \begin{center}
    \includegraphics[width=0.5\linewidth]{res/uni.png}
    \vspace{2cm}

    \huge{\textbf{Berechnung von Molekül-Grundzustandsenergien 
    mit Ab-Initio-Methoden}}
    \vspace{1cm}

    \Large
    Fachbereich Informatik

    Bachelorarbeit

    Deniz Güven
    \vfill

    Betreuer: 

    \today
  \end{center}
\end{titlepage}

\tableofcontents

%=========================================================================
\chapter{Einleitung}
\section{Chemischer Hintergrund}

TODO
Definition von Chemie

\subsection{Atommodelle}
In der Chemie wurden immer präzisere 
und umfangreichere Beschreibungen des Atoms entwickelt,
zu diesen gehören unter anderen:
TODO
\begin{enumerate}
  \item Demokrit
  \item Dalton
  \item Rutherford
  \item Bohr
  \item Quantenmechanische Modell
\end{enumerate}

\section{Ziele dieser Bachelorarbeit}
Eine wichtige Eigenschaft von Atomen und Molekülen ist die Grundzustandsenergie, 
mit dieser können viele Prozesse in der Chemie,
wie Reaktionsabläufe und Molekülstrukturen, erklärt werden.
Diese Bachelorarbeit befasst sich mit der Berechnung 
dieser Energie von einfachen Atomen und Molekülen innerhalb einer quantenmechanischen Beschreibung.
Dafür wird ein theoretisches Verständis entwickelt 
und diese Theorie dann in einem Programm implementiert.

\section{Relevanz für die Chemie}
TODO Genaue Verwendung von GZs erklären. 

%=========================================================================
\chapter{Theorie und Methoden}
\section{Allgemeine Theorie}
\subsection{Grundlegende Defininitionen}
\subsection{Postulate der Quantenmechanik}
TODO: $\psi$ als immer normiert annehmen?, Mehr infos + umschreiben?
\subsubsection{Wellenfunktionen}
Die sogenannte Wellenfunktion $\Psi(x_1, x_2, \dots, t)$ 
beschreibt den Zustand eines quantenmechanischen Systems vollständig.\cite[S. 20-21]{atkins_friedman_2011}\\\\
Durch ein $\psi$ wird nur der räumliche Anteil dargestellt.

\subsubsection{Operatoren}
Beobachtbare Eigenschaften eines quantenmechanischen Systems 
werden durch sogenannte Operatoren repräsentiert.
Diese Operatoren müssen hermitisch sein und folgende Gleichungen erfüllen:
\begin{flalign*}
  \hat{q}\hat{p}_{q'} - \hat{p}_{q'}\hat{q} &= i\hbar\delta_{qq'}\\
  \hat{q}\hat{q}' - \hat{q}'\hat{q} &= 0 \\
  \hat{p}_{q}\hat{p}_{q'} - \hat{p}_{q'}\hat{p}_{q} &= 0
\end{flalign*}
Wobei $\hat{q}, \hat{q}' \in \{\hat{x}, \hat{y}, \hat{z}\}$ die Orts-Operatoren stellen
und $\hat{p}_q, \hat{p}_{q'}$ die zugehörigen Impuls-Operatoren sind.\cite[S. 21]{atkins_friedman_2011}

\subsubsection{Messungen}
Der Mittelwert von wiederholten Messungen
entspricht dem Erwartungswert des korrespondierenden Operators auf der Wellenfunktion.\\
Der Erwartungswert eines Operators $\hat{o}$ ist gegeben durch:
\begin{equation}
  \left\langle \hat{o} \right\rangle = \frac{\int \psi^* \hat{o} \psi \,d\tau}{\int \psi^* \psi \,d\tau}
  =\frac{\left\langle \psi^* \vert \hat{o} \vert \psi \right\rangle}{\left\langle \psi^* \vert \psi \right\rangle}
\end{equation}\\
Wenn $\psi$ eine Eigenfunktion eines Operators $\hat{o}$ ist, erhält man den Eigenwert als Erwartungswert.
\cite[S. 22-23]{atkins_friedman_2011} 

\subsubsection{Bornsche Wahrscheinlichkeitsinterpretation}
Die Wahrscheinlichkeit ein Teilchen in einem Volumenelement $\,d\tau$ zu finden ist geich
$\vert \psi(x) \vert^2$, wenn $\psi$ normiert ist.\cite[S. 24]{atkins_friedman_2011}

\subsubsection{Die Schrödingergleichung}
Die zeitliche Änderung dieser Wellenfunktion $\Psi(x_1, x_2, \dots, t)$ 
wird durch die Schrödingergleichung beschrieben:
\begin{equation}
  i\hbar\frac{\partial\Psi}{\partial t} = \hat{H}\Psi
\end{equation}\\
Sollte die Zeitabhägigkeit der Wellenfunktion trivial sein (TODO erklären?),
kann diese in 2 Funktionen zerlegt werden:
\begin{equation*}
  \Psi(x_1, x_2, \dots, t) = \psi(x_1, x_2, \dots) \exp(\frac{-iEt}{\hbar})
\end{equation*}\\
Dabei ist $\psi$ die Lösung für die zeitunabhängige Schrödingergleichung:
\begin{equation}
  \hat{H}\psi = E\psi
\end{equation}
\cite[S. 24-25]{atkins_friedman_2011}

\subsection{Hamilton-Operator}
Der allgemeine Hamilton-Operator für Möleküle mit $N$ Elektronen und $M$ Atomkernen lautet:
\begin{flalign}\label{hamilton}
  \hat{H} &= -\sum_i^N \frac{1}{2} \nabla_i^2 
            -\sum_\mu^M \frac{1}{2 m_\mu} \nabla_\mu^2
            -\sum_i^N \sum_\mu^M \frac{Z_\mu}{r_{i\mu}}
            +\frac{1}{2}\sum_{i\neq j} \frac{1}{r_{ij}}
            +\frac{1}{2}\sum_{\mu \neq \nu } \frac{Z_\mu Z_\nu}{r_{\mu\nu}}\nonumber\\
          &= \hat{T}_e + \hat{T}_A + \hat{V}_{eA} + \hat{V}_{ee} + \hat{V}_{AA}
\end{flalign}
Dabei steht $Z_\mu$ für die Ladung des Atomkerns $\mu$ und 
$r$ für den Abstand zwischen Elektronen und/oder Atomkernen.
\cite[S. 6]{tc2_1}

\subsection{Born-Oppenheimer-Näherung}

Aufgrund des hohen Masseunterschiedes zwischen Elektronen und Atomkernen
ist der Einfluss der Elektronen auf die Bewegung der trägeren Atomkerne vernachlässigbar.
Deshalb werden die Atomkerne innerhalb der Born-Oppenheimer-Näherung als statisch betrachtet 
bei der Berechnung der Elektronen-Wellenfunktion.

Dafür wird der Allgemeine Hamilton-Operator \cref{hamilton} aufgeteilt:
\begin{flalign*}
  \hat{H} &= \hat{T}_e + \hat{T}_A + \hat{V}_{eA} + \hat{V}_{ee} + \hat{V}_{AA}\\
          &= \hat{T}_A + \hat{H}_{\text{el}}
\end{flalign*}
Wir lösen nun die Schrödingergleichung mit dem elektronischen Hamilton-Operator $\hat{H}_{\text{el}}$.
Bei der Lösung dieser wird eine feste Kerngeometrie angenommen, 
die wir bei den Operatoren $\hat{V}_{eA}$ und $\hat{V}_{AA}$ verwenden werden. 
Mit der resultierenden Elektronen-Wellenfunktionen 
lässt sich dann eine Gesamte Wellenfunktion mit Einbezug der Atomkerne konstruieren.
\cite[S. 11-14]{tc2_1}

\subsection{Variationsformulierung}
Da eine analytische Lösung zur Schrödingergleichung nur in speziellen Fällen existiert \cite[S. 195]{lewars_2016},
wird die exakte Wellenfunktion durch eine Test-Wellenfunktion approximiert.
Es lässt sich zeigen, 
dass die Energie dieser Test-Wellenfunktion $E_{test}$ 
immer über der tatsächlichen Grundzustandsenergie $E_0$ liegt.

\subsubsection*{Beweis}
TODO das N korrekt?\\
TODO Auf Rayleigh-ritz erweitern und auf det$|\hat{H}_{ij} - ES_{ij}| = 0$ kommen\\
TODO test funktion normiert annahme -> $E_0$ kann dann locker ins integral.
Voraussetzungen:
\begin{enumerate}
  \item Der Hamiltonian $\hat{H}$ hat die Eigenfunktionen $\psi_i^{}$ mit Eigenwerten $E_i^{}$.
  \item Es existiert eine Eigenfunktion $\psi_0^{}$ mit dem niedrigsten  Eigenwert $E_0^{}$.
  \item Alle Eigenfunktion von $\hat{H}$ sind orthonomal zueinander:\\
  $\forall i,j \in \mathbb{N}_0 : \langle \psi_i^{} \vert \psi_j^{} \rangle = \delta_{ij}^{}$
  \item Die Test-Wellenfunktion lässt sich als Lineakombination der Eigenfunktionen darstellen:
  $\psi_{test}^{} = \sum_{n}^{} c_n^{} \psi_n^{}$
\end{enumerate}
Zu zeigen: $E_{test}^{} \geq E_0^{}$ oder $E_{test}^{} - E_0^{} \geq 0$
\begin{flalign*}
  E_{test} - E_0 
  &= \langle \psi_{test}^{} \vert \hat{H} - E_0^{} \vert \psi_{test}^{} \rangle\\
  &= \int \psi_{test}^* (\hat{H} - E_0^{}) \psi_{test}^{} \,dx \quad &\vert \text{ 4. Voraussetzung}\\
  &= \sum_n \sum_m c_n^\ast c_m^{} \int \psi_{n}^* (\hat{H} - E_0^{}) \psi_{m} \,dx 
  \quad &\vert \text{ 1. Voraussetzung}\\
  &= \sum_n \sum_m c_n^\ast c_m^{} \int \psi_{n}^* (E_m^{} - E_0^{}) \psi_{m} \,dx 
  \quad &\vert \text{ 3. Voraussetzung}\\
  &= \sum_n \left\lvert c_n^{} \right\rvert^2 (E_n^{} - E_0^{}) \int \psi_{n}^* \psi_{n} \,dx 
  \quad &\vert \text{ 3. Voraussetzung}\\
  &= \sum_n \left\lvert c_n^{} \right\rvert^2 (E_n^{} - E_0^{})
  \quad &\vert \text{ 2. Voraussetzung}\\
  &\geq 0 &\qed
\end{flalign*}\\
\cite[S. 187]{atkins_friedman_2011}\\
%\cite[]{TC2}
TODO Warum Voraussetzungen erfüllt sind, erklären.

\subsection{Beschreibung von Elektronen}
Die gesamte Wellenfunktion $\Psi$ muss noch einige Eigenschaften erfüllen, 
um die Elektronen des Moleküls zu beschreiben:

\begin{enumerate}
  \item Jedes Elektron verfügt, neben der räumlichen Ausdehnung(TODO: Formulierung), 
  auch über einen intrinsischen Spin. Dieser Spin kann zwei Zustände annehmen, diese werden über
  die Spinfunktionen $\alpha$ und $\beta$ dargestellt.
  \item Die Gesamt-Wellenfunktion muss antisymmetrisch sein 
  bezüglich dem Austausch von zwei beliebigen Elektronen z.B.
  $\Psi(1, 2) = - \Psi(2, 1)$
\end{enumerate}
TODO cite\\
Deshalb stellen wir die Gesamtwellenfunktion mit $2n, n \in \mathbb{N}$ Elektronen als eine Determinante von Spin-Orbitalfunktionen dar:
\begin{flalign*}
&\Psi(1, 2, \dots, n) =\\ 
&\frac{1}{\sqrt{(2n)!}}\left\lvert
\begin{array}{ccccc} 
\psi_1(1)\alpha(1) & \psi_1(1)\beta(1) & \cdots & \psi_n(1)\alpha(1) & \psi_n(1)\beta(1)\\ 
\psi_1(2)\alpha(2) & \psi_1(2)\beta(2) & \cdots & \psi_n(2)\alpha(2) & \psi_n(2)\beta(2)\\ 
    \vdots         &       \vdots      & \ddots &       \vdots       &       \vdots     \\ 
\psi_1(2n)\alpha(2n) & \psi_1(2n)\beta(2n) & \cdots & \psi_n(2n)\alpha(2n) & \psi_n(2n)\beta(2n)\\ 
\end{array}\right\rvert
\end{flalign*}

TODO zusatzinfos + eqn number? Erklärung spin auf s.200 lewars


\section{Hartree-Fock}

Im Hartree-Fock-Ansatz wird aus dem Hamilton-Operator, der auf die gesamte Wellenfunktion wirkt, 
ein 1-Elektronen-Operator entwickelt, der Fock-Operator $\hat{F}$.
Da dieser Operator nur auf einzelne Elektronen wirkt, 
kann eine Pseudo-Eigenwerts-Gleichung für jedes Orbital erstellt werden.

Außerdem ist in diesem Ansatz die Elektron-Elektron-Abstoßung nur approximativ behandelt.
Dabei wird vereinfacht für die Repulsionen nur die Abstoßung die ein Elektron in einem gemittelten Potential
, das durch die anderen Elektronen entsteht, berechnet.
Für eine genauere Beschreibung müsste die Abstoßung aller Elektronen-Paare individuell berücksichtigt werden.
TODO Visualisierung.

Über die Hartree-Fock-Methode folgen diese Gleichung für jedes Spinorbital:
\begin{equation}
  \hat{F} \psi_i = \epsilon_i \psi_i \quad i = 0 \dots n
\end{equation}

\subsection{Fock Operator}


- Lösungsweg über das SCF-Verfahren (Matrix-Darstellung, ...)

- Verwendung von Basisfunktionen (Konstruktion der Wellenfunktion)

- Implementierung (größten Probleme: Integral-Evaluierung und
Matrix-Diagonalisierung)

\section{DFT}
- Herleitung (Nur die Idee/Ergebnisse, da wahrscheinlich über
meinem Niveau)

- Konkrete Umsetzung durch die Kohn-Sham-Gleichung (Terme in der
Schrödinger-Glg. + XC-Funktionale)

- Lösung durch FEM + PINVIT (+ LDA)

- Implemtierung über UG4 LUA

%=========================================================================
\chapter{Ergebnisse/ Numerische Experimente}
\section{Erklärung der Experimente}
- Eigen-Energien als Benchmark + Moleküle zum Testen (Simple wie
H2O, CH4, ... und Komplexe wie z.b. Benzol, das eine
Elektronen-Delokalisation aufweist)

\section{Experimente(HF, DFT, FULL-CI(exakt) über NWCHEM oder Literatur)}
- Präsentation der Ergebnisse(Graphen, Tabellen, usw.)
- Werden Effekte bei komplexen Molekülen korrekt erfasst?

\section{Vergleich der Methoden/Deutung der Ergebnisse (HF vs. DFT)}
- Genauigkeit, Kosten, Skalierbarkeit, ...

%=========================================================================
\chapter{Diskussion/Ausblick}
\section{Einordnung von HF und DFT in der Chemie}
-> Andere Klassen von Methoden (zb. semiempirische Methoden)
-> Verbesserung dieser Methoden (Post-Hartree-Fock-Methoden)
-> Eingliederung dieser Methoden in der Praxis (Was kann man mit
diesen Eigenenergien/Funktionen eigentlich machen?).

\section{Wie könnte man von diesem Punkt aus weitermachen?}
-> Code-Optimierung, Anspruchsvoller Methoden implementieren
(aufbauend auf HF), Geometrie-Optimierung, ...

%=========================================================================
\bibliographystyle{IEEEtran}
\bibliography{IEEEabrv, main}

\end{document}